% unidades/05-youto-suru.tex
\chapter{\emoji{running} \ruby{ようとする}{} — Intentar / Negarse}
\unidadheader{05}{意志}{ようとする / ようとしない}{Expresar intento de acción o resistencia}

\begin{tcolorbox}[vocabbox, title={\emoji{pushpin} Vocabulario}]
\begin{tabularx}{\linewidth}{llX}
\toprule
\textbf{Japonés} & \textbf{Español} & \textbf{Tipo} \\ \midrule
\vocabitem{逃げる}{にげる}{huir / escapar}{v. G2}
\vocabitem{起きる}{おきる}{levantarse / despertarse}{v. G2}
\vocabitem{諦める}{あきらめる}{rendirse / abandonar}{v. G2}
\vocabitem{挑戦する}{ちょうせんする}{desafiar / intentar}{v. suru}
\vocabitem{謝る}{あやまる}{disculparse}{v. G1}
\vocabitem{認める}{みとめる}{reconocer / admitir}{v. G2}
\vocabitem{拒否する}{きょひする}{rechazar / negarse}{v. suru}
\bottomrule
\end{tabularx}
\end{tcolorbox}

\subsection*{\emoji{speech-balloon} Frases Clave}
\frase{\ruby{起}{お}きようとしたが、できなかった。}{Intenté levantarme pero no pude.}
\frase{\ruby{逃}{に}げようとしても\ruby{無駄}{むだ}です。}{Es inútil intentar huir.}
\frase{\ruby{彼}{かれ}は\ruby{謝}{あやま}ろうとしない。}{Él se niega a disculparse.}
\frase{\ruby{諦}{あきら}めようとしたとき、\ruby{奇跡}{きせき}が\ruby{起}{お}きた。}{Justo cuando estaba a punto de rendirme, ocurrió un milagro.}

\subsection*{\emoji{gear} Gramática}
\patron{ようとする — estar a punto de / intentar}
  {[V-\ruby{意向形}{いこうけい}] + とする}
  {\ej{\ruby{寝}{ね}ようとしたとき、\ruby{電話}{でんわ}が\ruby{鳴}{な}った。}{Justo cuando estaba a punto de dormir, sonó el teléfono.}
   \ej{ドアを\ruby{開}{あ}けようとしたが、\ruby{鍵}{かぎ}がかかっていた。}{Intenté abrir la puerta, pero estaba con llave.}
   \ej{\ruby{立}{た}ち\ruby{上}{あ}がろうとしている。}{Está intentando ponerse de pie.}}

\patron{ようとしない — negarse a / no querer}
  {[V-\ruby{意向形}{いこうけい}] + としない}
  {\ej{\ruby{彼}{かれ}は\ruby{謝}{あやま}ろうとしない。}{Él se niega a disculparse.}
   \ej{\ruby{子供}{こども}が\ruby{野菜}{やさい}を\ruby{食}{た}べようとしない。}{El niño se niega a comer verduras.}}

\comparacion{ようとする vs ようとしない}
  {ようとする\\\small intento / a punto de\\\small \ruby{起}{お}きようとした\\\textit{Intentó levantarse}}
  {ようとしない\\\small resistencia / negativa\\\small \ruby{起}{お}きようとしない\\\textit{Se niega a levantarse}}
  {La misma base (forma volitiva + と) con する o しない cambia completamente el significado: esfuerzo vs resistencia.}

\coloquial{ようとしている}{ようとしてる / ようとしてんの？}{
  En habla casual: としている → としてる。 Para preguntar con incredulidad: \ruby{逃}{に}げようとしてんの？ ¿Estás intentando escapar?}

\culturabox{\ruby{意志}{いし}の\ruby{表現}{ひょうげん} — Expresar la voluntad en japonés}{
  La forma volitiva (\ruby{意向形}{いこうけい}) expresa la voluntad del hablante. Al combinarse con とする, el foco pasa del resultado a la intención. En la cultura japonesa, expresar que se «intenta» algo (esfuerzo) es valorado incluso cuando no se logra: \ruby{頑張}{がんば}ろうとした (intenté esforzarme).
}

\lectura{\ruby{彼女}{かのじょ}は\ruby{何度}{なんど}も\ruby{立}{た}ち\ruby{上}{あ}がろうとした。\ruby{足}{あし}が\ruby{痛}{いた}くて、なかなかできなかった。それでも\ruby{諦}{あきら}めようとしなかった。\ruby{周}{まわ}りの\ruby{人}{ひと}が\ruby{応援}{おうえん}してくれて、ついに\ruby{成功}{せいこう}した。}
{Ella intentó ponerse de pie muchas veces. Le dolían los pies y no lo lograba fácilmente. Aun así, no quiso rendirse. Las personas a su alrededor la animaron y al final lo logró.}

\begin{tcolorbox}[ejerciciobox, title={\emoji{pencil} Ejercicios}]
\begin{enumerate}
  \item Conjuga en forma volitiva: \ruby{食}{た}べる、する、\ruby{来}{く}る、\ruby{書}{か}く
  \item Transforma: \ruby{彼}{かれ}は\ruby{認}{みと}めない (→ usa ようとしない)
  \item Crea una oración con ようとしたが usando una experiencia propia.
\end{enumerate}
\textbf{Respuestas:}
1) \ruby{食}{た}べよう・しよう・こよう・\ruby{書}{か}こう \quad
2) \ruby{彼}{かれ}は\ruby{認}{みと}めようとしない \quad
3) (libre)
\end{tcolorbox}

\begin{tcolorbox}[kanjibox, title={\emoji{mahjong} Kanjis}]
\begin{tabularx}{\linewidth}{cllXX}
\toprule
\textbf{Kanji} & \textbf{On} & \textbf{Kun} & \textbf{Significado} & \textbf{Vocab} \\ \midrule
\kanjirow{逃}{とう}{\ruby{に}{げる}}{huir}{\ruby{逃}{に}げる / \ruby{逃亡}{とうぼう}}
\kanjirow{挑}{ちょう}{\ruby{いど}{む}}{desafiar}{\ruby{挑戦}{ちょうせん} / \ruby{挑}{いど}む}
\kanjirow{諦}{てい}{\ruby{あきら}{める}}{resignarse}{\ruby{諦}{あきら}める}
\kanjirow{認}{にん}{\ruby{みと}{める}}{reconocer}{\ruby{認}{みと}める / \ruby{認識}{にんしき}}
\bottomrule
\end{tabularx}
\end{tcolorbox}
